\documentclass[11pt,a4paper]{article}

\usepackage[T1]{fontenc}
\usepackage[utf8]{inputenc}
\usepackage[french]{babel}
\usepackage{lmodern}
\usepackage{microtype}
\usepackage{geometry}
\usepackage{graphicx}
\usepackage{booktabs}
\usepackage{amsmath}
\usepackage{amssymb}
\usepackage{hyperref}

\geometry{margin=2.5cm}

\title{Titre de l'article}
\author{Prenom Nom\\Institution\\\texttt{email@example.com}}
\date{\today}

\begin{document}

\maketitle

\begin{abstract}
Cet abstract resume le contexte, la methode, les resultats principaux et la contribution de l'article en quelques phrases.
\end{abstract}

\section{Introduction}

Presenter le probleme scientifique, le contexte bibliographique et la contribution principale de l'article.

\section{Methode}

Decrire les donnees, les hypotheses et la methode utilisee.

\begin{equation}
  y = f(x) + \varepsilon
\end{equation}

\section{Resultats}

Presenter les resultats de maniere concise.

\begin{table}[h]
  \centering
  \caption{Exemple de tableau de resultats.}
  \begin{tabular}{lcc}
    \toprule
    Methode & Precision & Rappel \\
    \midrule
    Reference & 0.82 & 0.78 \\
    Proposee & 0.90 & 0.86 \\
    \bottomrule
  \end{tabular}
  \label{tab:resultats}
\end{table}

\section{Discussion}

Interpretez les resultats, leurs limites et leur portee.

\section{Conclusion}

Rappeler la contribution principale et ouvrir sur les perspectives.

\begin{thebibliography}{9}

\bibitem{exemple}
Auteur, A.
\newblock Titre de la reference.
\newblock \emph{Journal}, 1(1), 1--10, 2026.

\end{thebibliography}

\end{document}

